\documentclass[12pt]{article}

\usepackage[a4paper, total={6in, 8in}]{geometry}
\usepackage{textcomp}

\begin{document}
\setlength{\parindent}{0pt}

\def \t {\textrightarrow}
\def \v {\vspace{1em}}
\def \bi {\begin{itemize}}
\def \ei {\end{itemize}}
\def \s[#1] {\section*{#1}}
\def \ss[#1] {\subsection*{#1}}
\def \sss[#1] {\subsubsection*{#1}}

\s[Fosfogliceridi]
Hanno due acidi grassi, e poi un fosfato legato un altra catena (di solito è un amminoalcol) \\
Estereficazione si crea il legame COO \t COH + OH \t condensazione, si elimina l'aqua \\
Acido fosfatifico \t importante nei legami forti, si vede anche nell ATP \\
Testa idrofila \t carica negativa ? che porta a creazione di legami a idrogeno \\
Teste polari sono a contatto con il liquido \t mentre le code idrofobiche all opposto

\s[Terpeni]
2-metil-1,3-butandiene (diene perchè ha due doppi legami) \t questo è l'isoprene \\
È una molecola instabile \t tende alla polimerizzazione

\s[Colesterolo]
I saponi sono dei sali a catena lunga \\
I terpeni non sono saponificabili \t perchè non c'è la parte che puo formare il sale: isoprene non ha acidi grassi, non ha la parte che puo formare il sale \\
Anche il colesterolo \t 3 esagoni e un pentagono \t è uno steoride = gruppo di lipidi saponoficabili con idrucarburo policiclico saturo (saturo ovvero presenta doppio legame) \\
Alchile = catena dove i carboni sono legati con legami semplice \\
Colesterolo lo assumiamo ma lo produciamo anche \t e la produzione si regola con gli ormoni \t ma se lo introduco con alimentazione, non lo controllo \\
Colesterolo determina fluidita della membrana cellulare \t e ogni cellula ha sua funzionalita per cui serve una diversa fluiditia \t il colesterolo regola questo \\
Reazioni anabolizzanti = reazioni che costruiscono \t quando digerisco invece è reazione catabolizzatrice (distruggo) \\
Quando ho bisogno di colesterolo, chiedo al fegato di anabolizzare \\
Inoltre colesterolo impedisce che la membrana cristallizzi

\v

Sali biliari \t bile serve ad emulsionare i grassi = fenomeno per cui i lipidi vengono separati in piccole porzioni che vengono isolate tra loro da queste molecole, e questo rende + trasportabili e digeribili

\v

Ormoni sono idrofili \t mentre ormoni tiroidei e quelli delle gonadi sono lipofili \\
Cortisone ha una funzione antifiammatoria
\end{document}
